\def\EditionMode{0}
\def\EditionLabel{Isolated diplomatic proof}
\def\EditionStatus{Local layout verification}
\documentclass[10pt,oneside]{book}
\usepackage{fontspec}
\setmainfont{Latin Modern Roman}
\setsansfont{Latin Modern Sans}
\usepackage{amsmath,amssymb,mathtools}
\usepackage{unicode-math}
\setmathfont{Latin Modern Math}
\usepackage{xcolor,tikz-cd,lua-ul,marginnote,pdflscape}
\usetikzlibrary{arrows.meta,decorations.pathmorphing,calc,shapes.geometric}
\tikzcdset{ampersand replacement=\&}
\tikzcdset{every label/.append style={text=black}}
\usepackage[a4paper,left=18mm,right=45mm,top=18mm,bottom=22mm,marginparwidth=28mm,marginparsep=7mm]{geometry}
\usepackage[unicode,colorlinks=true,urlcolor=black,linkcolor=black]{hyperref}
\usepackage{bookmark}
\setlength{\parindent}{0pt}
\setlength{\parskip}{4pt plus 1pt minus 1pt}
\setlength{\emergencystretch}{1.5em}
\raggedbottom
\definecolor{PSHand}{RGB}{25,58,132}
\definecolor{PSDelete}{RGB}{146,37,37}
\definecolor{PSTypedOld}{RGB}{95,95,95}
\ifnum\EditionMode=2\colorlet{blue}{black}\else\colorlet{blue}{PSHand}\fi
\RenewDocumentCommand{\textcircled}{m}{\tikz[baseline=(PSCircle.base)]\node[draw,circle,inner sep=.5pt,line width=.4pt](PSCircle){#1};}
\newif\ifPSInApparatus
\newcounter{psnote}
\newcounter{psvariant}
\newcounter{pslinkmark}
\newif\ifPSInLink
\newsavebox{\PSNoteBox}
\newdimen\PSNoteLimit
{\footnotesize\global\PSNoteLimit=3\baselineskip}
\newcommand{\CurrentScanURL}{}
\newcommand{\CurrentKey}{}
\newcommand{\CurrentScan}{}
\newcommand{\pssource}[3]{\psflushmargins\clearpage\gdef\CurrentKey{#1}\gdef\CurrentScan{#2}\gdef\CurrentScanURL{#3}\setcounter{psvariant}{0}\typeout{PS-SOURCE: #1}}
\newcommand{\psdivision}[2]{\pdfbookmark[0]{#1}{#2}\par\noindent{\large\bfseries #1\par}\medskip}
\newcommand{\pslinkmarker}[2]{\stepcounter{pslinkmark}\latelua{texio.write_nl('log','PS-LINK-SHIP|#1|\CurrentKey|\arabic{pslinkmark}|#2')}}
\newcommand{\pslink}[2]{\ifPSInLink\pslinkmarker{#1}{nested}#2\else\begingroup\PSInLinktrue\href{\CurrentScanURL}{\pslinkmarker{#1}{outer}#2}\endgroup\fi}
\newcommand{\scanpage}[1]{\ifnum\EditionMode=2\leavevmode\marginpar{\scriptsize\sffamily\pslink{scanpage}{scan #1}}\else\par\noindent{\small\sffamily\pslink{scanpage}{[scan #1]}}\par\medskip\fi}
\newcommand{\mspage}[1]{\ifnum\EditionMode=2\leavevmode\marginpar{\scriptsize\sffamily p.\,#1}\else{\small\sffamily[p.\,#1]}\fi}
\newcommand{\hand}[1]{{\ifnum\EditionMode=2\color{black}\else\color{PSHand}\fi #1}}
\let\over\relax
\NewDocumentCommand{\over}{O{} m m}{{\textcolor{PSTypedOld}{#2}\ensuremath{^{\!\nearrow}}\if\relax\detokenize{#1}\relax\hand{#3}\else #3\fi}}
\NewDocumentCommand{\inserted}{O{} m}{{\if\relax\detokenize{#1}\relax\hand{\ensuremath{\langle}#2\ensuremath{\rangle}}\else\ensuremath{\langle}#2\ensuremath{\rangle}\fi}}
\newcommand{\deleted}[1]{{\color{PSDelete}\ifmmode\text{[del.: }#1\text{]}\else[\textsf{del.:} #1]\fi}}
\newcommand{\xout}[1]{{\color{PSTypedOld}\ifmmode\text{[x: }#1\text{]}\else[\textsf{x:} #1]\fi}}
\let\PSMathUnderline\underline
\RenewDocumentCommand{\underline}{+m}{\ifmmode\PSMathUnderline{#1}\else\underLine{#1}\fi}
\NewDocumentCommand{\uline}{+m}{\underline{#1}}
\NewDocumentCommand{\uuline}{+m}{\ifmmode\PSMathUnderline{\PSMathUnderline{#1}}\else\underLine{\underLine{#1}}\fi}
\newunderlinetype\PSWaveLine{\leaders\hbox{\lower.7ex\hbox{\tikz[x=1pt,y=1pt,baseline=0pt]{\draw[line width=.3pt] (0,0) .. controls (1,1) and (2,1) .. (3,0) .. controls (4,-1) and (5,-1) .. (6,0);}}}}
\NewDocumentCommand{\uwave}{+m}{{\PSWaveLine #1}}
\newcommand{\unsure}[1]{\pslink{unsure}{[#1\ifmmode\,?\else\,?\fi]}}
\newcommand{\illegible}[1]{\ifmmode\text{\pslink{illegible}{[illeg.\ #1]}}\else\pslink{illegible}{[illeg.\ #1]}\fi}
\NewDocumentCommand{\diagram}{+m}{\ifmmode\vcenter{\hbox{#1}}\else\par\medskip{\centering #1\par}\medskip\psflushmargins\fi}
\newenvironment{diagramblock}{\par\begin{center}}{\end{center}\par}
\newcommand{\PSMathNotes}{}
\makeatletter
\NewDocumentCommand{\psplacednote}{m +m}{%
      \sbox{\PSNoteBox}{\parbox{\marginparwidth}{\footnotesize\raggedright\hand{\textsf{m#1.} #2}}}%
      \ifdim\dimexpr\ht\PSNoteBox+\dp\PSNoteBox\relax<\PSNoteLimit
        \leavevmode\marginpar{\usebox{\PSNoteBox}}%
      \else
        \par\nobreak\smallskip\hrule\nobreak\smallskip
        {\footnotesize\textsf{Marginal note m#1.}\quad\hand{#2}\par}
        \nobreak\smallskip\hrule\smallskip
      \fi}
\newcommand{\psflushmargins}{\ifmmode\else\ifinner\else\ifx\PSMathNotes\@empty\else\let\PSTempNotes\PSMathNotes\global\let\PSMathNotes\@empty\PSTempNotes\fi\fi\fi}
\NewDocumentCommand{\margin}{+m}{%
  \ifPSInApparatus\ifmmode\hand{\text{[marg.: }#1\text{]}}\else\hand{\textsf{[marg.: }#1\textsf{]}}\fi
  \else\stepcounter{psnote}%
    \ifmmode
      \text{\textsuperscript{m\arabic{psnote}}}%
      \begingroup\edef\PSQueue{\endgroup\noexpand\g@addto@macro\noexpand\PSMathNotes{\noexpand\psplacednote{\arabic{psnote}}{\unexpanded{$#1$}}}}\PSQueue
    \else
      \textsuperscript{m\arabic{psnote}}\psplacednote{\arabic{psnote}}{#1}%
    \fi
  \fi}
\AtBeginDocument{\everymath\expandafter{\the\everymath\aftergroup\psflushmargins}\everydisplay\expandafter{\the\everydisplay\aftergroup\psflushmargins}}
\makeatother
\newcommand{\psnumber}[2]{\expandafter\gdef\csname psnum@#1\endcsname{#2}}
\NewDocumentCommand{\var}{+m +m m m}{\ifnum\EditionMode=0 #1\else\ifnum\EditionMode=1\ifmmode\text{\textsuperscript{\tiny\csname psnum@#4\endcsname}}\else\textsuperscript{\tiny\csname psnum@#4\endcsname}\fi\fi #2\fi}
\newcommand{\psapparatusstart}{\par\bigskip\hrule\smallskip\PSInApparatustrue\begingroup\footnotesize\raggedright\textsf{Apparatus for scan \CurrentScan}\par}
\newcommand{\psapparatusend}{\par\endgroup\PSInApparatusfalse}
\NewDocumentCommand{\psentry}{m m m +m +m m}{\par\noindent\textbf{#2. #3.}\quad\textit{Typescript:} #4\quad\textit{Corrected:} \if\relax\detokenize{#5}\relax\textnormal{[omitted]}\else #5\fi\quad #6\par}
\newcommand{\psfrontmatter}{%
\hypersetup{pageanchor=false}\begin{titlepage}\raggedright
{\Large Alexander Grothendieck\par}\vspace{12mm}
{\Huge Pursuing Stacks\par}\medskip{\LARGE 1983\par}\vspace{12mm}
{\Large\EditionLabel\par}\vspace{8mm}
{\large\EditionStatus\par}\vfill
Université de Montpellier, Grothendieck archive, cotes 134-2--134-8.\par
Including the 26 source pages of R. Brown's translation at cote 134-2, sites 115--143.\par
\vspace*{2mm}\end{titlepage}\hypersetup{pageanchor=true}
\chapter*{Editorial statement}
AI-transcribed by Claude Opus 5.5 (Anthropic, effort high); checked against the scans, and this edition built, by GPT-6 Astra (OpenAI, ChatGPT Pro mode); prompts and records by Claude (Anthropic).

The witness consulted and not reproduced is the LaTeX edition of Mateo Carmona with the collaboration of Ulrik Buchholtz (arXiv 2111.01000 v2, 2021), extended from the 2015 edition by ``the scrivener'', CC0 1.0.

The translation in the appendix is R. Brown's, prepared with assistance from L. Breen and J-L. Loday, as his note says; the rights in it belong to its translator or his heirs.

The rights in the original work belong to its author or the author's representatives or heirs. Where we hold rights, in the AI-produced transcription, translation, corrections and typesetting, they are dedicated to the public domain under CC0 1.0. This edition is published for educational, non-profit purposes, and we claim no rights we do not hold. If you hold rights in this work and object to its publication here, please open a GitHub issue and we will take it down.

\section*{How to read this form}
\ifnum\EditionMode=0
The diplomatic form preserves typed and handwritten wording, cancellations, insertions, overwritten readings, underlining, doubts, and original page numbers. Blue marks handwriting; red labels handwritten cancellations; gray labels typewriter cancellations and underlying overwritten text. Angle brackets mark insertions; a superscript arrow separates an overwritten reading from its replacement. Typed insertions and typed overwriting remain plain. The disposition of long marginal notes is identified explicitly.
\else
The corrected text applies the author's insertions, replacements, cancellations, and recorded transpositions. It retains his English and notation, including consistent forms such as ``standart'', ``developped'', ``it's'' for ``its'', and ``A la''. An isolated slip is corrected only when specifically recorded and justified. Mathematical anomalies are preserved and discussed in the guide.
\ifnum\EditionMode=1
Blue continues to identify handwriting. A numbered call identifies each maximal editorial operation, with its full diplomatic and corrected readings and reason in the apparatus following the source page. Internal layers of a wholly replaced or deleted passage are preserved within that entry rather than duplicated as separate entries. These entries are generated from the same variants that govern the running text.
\else
This reading form omits the apparatus, cancelled material, superseded readings, insertion brackets, and color distinctions. It retains the resulting text, marginal notes, underlining, and unresolved doubts. Scan and source-page references appear in the margin.
\fi\fi

Every scan mark and remaining doubt links to its source page in the archive viewer. Doubts are editorial reading marks, not emendations. The source-page breaks identify the archival sequence; a source page may occupy more than one printed page. A partial working edition is explicitly labelled and contains only reviewed source pages.
\clearpage}

\begin{document}
\pssource{ps_134-6_067}{134-6/67}{https://grothendieck.umontpellier.fr/web/viewer.html?file=../134-6.pdf\#page=68}
% Pursuing Stacks (cotes 134-2 to 134-8): archive cote 134-6, site page 67 (PDF page 68), typescript page 385.
% typescript. Draft compared: ps5.tex:4756-4839.
% AI-transcribed by Claude Opus 5.5 (claude-opus-5-5, adaptive thinking, effort high), Message Batches API, batch ps3 (msgbatch_01XdFmuVVBAcgc7EFxdtCqsj), 2026-11-12; reviewed against both supplied scan halves; see CHECK_DECISIONS.tsv.
\scanpage{134-6/67}
\mspage{385}

For the sake of brevity, we'll say that a \uline{U}-category satisfying the
assumptions of prop.1 is a \hand{\uline{\textcolor{black}{pseudo}}}-\hand{\uline{\textcolor{black}{topos}}} (as these conditions are satisfied
for any topos). We get at once the

\hand{\uline{\textcolor{black}{Corollary}}} 2. Let \hand{$\mathcal{A}$},\hand{$\mathcal{B}$} be two pseudo-topoi. Then a functor from
\hand{$\mathcal{A}$} to \hand{$\mathcal{B}$}$^{o}$ (resp. from \hand{$\mathcal{B}$}$^{o}$ to \hand{$\mathcal{A}$} ) has a right adjoint (resp. a left
adjoint) iff it commutes to small direct limits (resp. to small inverse
limits). Thus taking right and left adjoints we get two equivalences of
categories, quasi-inverse to each other
\[ (4) \qquad \underline{Hom}_{\hand{!}}(\hand{\mathcal{A}},\hand{\mathcal{B}}^{o})\hand{{}^{o}}\hand{\leftrightarrows}\underline{Hom}^{\hand{!}}(\hand{\mathcal{B}}^{o},\hand{\mathcal{A}})\deleted{\hand{{}^{o}}} \quad , \]
and the two members of (4) are canonically equivalent to the \inserted[typed]{\deleted{dual of the}} category
\[ (5) \qquad \underline{Hom}^{\overset{\hand{\illegible{1}}}{\hand{!!}}}_{\deleted{\illegible{1}}}(\hand{\mathcal{A}^{o}},\hand{\mathcal{B}^{o}};\text{ (Sets)}) \]
of functors
\[ \hand{\mathcal{A}^{o}}\text{x}\hand{\mathcal{B}^{o}\to}\text{(Sets)} \]
which commute with small \xout{direct}\inserted[typed]{inverse} limits with respect to either variable
(the other being fixed), \inserted[typed]{(5) being} viewed as a full subcategory of \uline{Hom}(\hand{$\mathcal{A}^{o}$}x\hand{$\mathcal{B}^{o}$},(Sets\hand{)})

\xout{Thex\allowbreak{}bifunctorx\allowbreak{}x\allowbreak{}x\allowbreak{}x\allowbreak{}$\overset{o}{\text{x\allowbreak{}}}$Bx\allowbreak{}x\allowbreak{}x\allowbreak{}x\allowbreak{}x\allowbreak{}(Sets)x\allowbreak{}x\allowbreak{}associatedx\allowbreak{}tox\allowbreak{}ax\allowbreak{}functor}
\hand{\uline{\textcolor{black}{Remarks}}}. \hand{1)}The first hand side of (4)\deleted{$_{o}$} is tautologically isomorphic to
the category Hom$^{\hand{!}}$(\hand{$\mathcal{A}$}$^{o}$,\deleted{\illegible{1}}\hand{$\mathcal{B}$}) (as for any two\inserted[typed]{U-}categories we have the
\xout{is of course} tautological \deleted{equivalence} \hand{isomorphism}
\[ (6) \qquad \deleted{(x,y)xx} \quad (\underline{Hom}_{\hand{!}}(\hand{\mathcal{P}},\hand{\mathcal{Q}}))^{o}\hand{\simeq}\underline{Hom}^{\hand{!}}(\hand{\mathcal{P}^{o}},\hand{\mathcal{Q}^{o}}) \quad ) \quad , \]
thus the equivalence (4) can be seen more symetrically as an equivalence
\[ (4') \qquad \underline{Hom}^{\hand{!}}(\hand{\mathcal{A}}^{o},\hand{\mathcal{B}})\hand{\approxeq}\underline{Hom}^{\hand{!}}(\hand{\mathcal{B}}^{o},\hand{\mathcal{A}}) \quad , \]
both categories being equivalent to (5) by using the equivalence (3)
for the second, and the corresponding equivalence for \hand{$\mathcal{B}$} for the first,
plus the tautological isomorphism
\[ (7) \qquad \underline{Hom}^{\hand{!}}(\hand{\mathcal{P}},\underline{Hom}^{\hand{!}}(\hand{\mathcal{Q}},\hand{\mathcal{M}})) \quad \underline{Hom}^{\hand{!!}}(\hand{\over{\illegible{1}}{\mathcal{P}}\deleted{\illegible{1}},\over{\illegible{1}}{\mathcal{Q}};\over{\illegible{1}}{\mathcal{M}}}) \quad , \]
for any three \uline{U}-categories \hand{$\mathcal{P}$},\hand{$\mathcal{Q}$},\hand{$\mathcal{M}$}.

\hand{2)} When \hand{$\mathcal{A}$} is a topos, \hand{$\mathcal{B}$} any \uline{U}-category stable under\inserted[typed]{small} inverse limits,
then we may interprete the category \uline{Hom}$^{\hand{!}}$(\hand{$\mathcal{A}$}$^{o}$,\hand{$\mathcal{B}$}) as the category of
\hand{$\mathcal{B}$}-\hand{\uline{\textcolor{black}{valued}}} \hand{\uline{\textcolor{black}{sheaves}}} on the topos (defined by)\hand{$\mathcal{A}$}. When \hand{$\mathcal{A}$} and \hand{$\mathcal{B}$} are both
topoi, then the equivalence (4) states that \xout{shxx} \hand{$\mathcal{B}$}-valued sheaves on
\hand{$\mathcal{A}$} can be identified with \hand{$\mathcal{A}$}-valued sheaves on \hand{$\mathcal{B}$}, and both may be
identified with set-valued ``bi-sheaves'' on \hand{$\mathcal{A}$} x \hand{$\mathcal{B}$}. In case the topoi
\hand{$\mathcal{A}$},\hand{$\mathcal{B}$} are defined respectively by\inserted{\underline{U}-}sites A,B (not necessarily small ones),

\margin{67}

\end{document}
